\section{Dataset Curation}\label{sec:data}

\subsection{Sampling}
We drew the repositories from SEART GitHub Search~\cite{dabic2021seart}, exported on 28~October~2026 with the filters in Table~\ref{tab:frame}.
Repositories created before 2022 existed through both events, and a commit after the last cutoff (February~2025) shows that a project was still maintained when a migration became unavoidable.
The thresholds on stars, contributors and commits restrict the frame to established, collaborative projects, the population for which questions about who migrates and how quickly are meaningful.
The export holds 22,608 repositories; we dropped the 393 archived ones, which can no longer change.
We shuffled the remaining 22,215 with a fixed seed and took the first 500, a simple random sample that later batches can extend.
One sampled repository no longer existed, and two sampled names turned out to be the same repository under its old and new name, which we counted once.
The sample therefore has 498 distinct repositories.

\begin{table}[t]
\caption{Sampling frame (SEART GitHub Search, 2026-10-28)}\label{tab:frame}
\centering\small
\begin{tabular}{@{}ll@{}}
\toprule
\textbf{Filter} & \textbf{Value} \\
\midrule
Created & before 2022-01-01 \\
Last commit & on or after 2025-02-01 \\
Stars / contributors / commits & $\geq$ 500 / $\geq$ 10 / $\geq$ 500 \\
Forks, archived & excluded \\
\midrule
Frame / sample & 22,215 / 498 repositories \\
\bottomrule
\end{tabular}
\end{table}

\subsection{Data collection}
We cloned every repository with its full commit history (without file contents, which git fetches on demand) on 28~October~2026, and analyzed the history of its default branch up to that day.
For each event and repository we scanned the files in \texttt{.github/workflows} at every relevant point of the first-parent history of the default branch, so that a change counts from the moment it reached the default branch, whether through a direct commit or a merge.

\textbf{Detection.}
For the artifact actions, a deprecated use is a \texttt{uses:} reference to \texttt{actions/upload-artifact} or \texttt{actions/download-artifact} at major version~1, 2 or~3, written as a major tag (\texttt{@v3}), a minor or patch tag (\texttt{@v3.1.3}), or a commit hash of one of those releases (we collected the hashes of all official v1--v3 tags).
Actions of other owners with similar names, such as \texttt{dawidd6/action-download-artifact}, do not count.
For the workflow commands, a deprecated use is the string \texttt{::set-output} or \texttt{::save-state}.
Commented-out lines are ignored for both.

\textbf{Exposure and migration.}
A repository is \emph{exposed} to an event if its default branch contained at least one deprecated use on the day of the deprecation notice ($T_{warn}$).
For an exposed repository, the \emph{migration} is the earliest default-branch commit after $T_{warn}$ after which no deprecated use remains; a repository with uses left on the collection day did not migrate.
For the artifact actions we treat v1/v2 and v3 separately, because their cutoffs differ: a migration counts as after the cutoff if any version group that the repository used at $T_{warn}$ disappeared only after that group's cutoff.

Of the 498 repositories, 124 were exposed to at least one event: 90 to the artifact deprecation, 57 to the \so{} deprecation and 23 to both, giving 147 (repository, event) pairs.

\subsection{Validation}
We developed the detection rules on a separate, hand-picked set of 90 popular repositories, checking every detected case against its diff and commit message.
This revealed the problems the final rules address: commented-out lines, look-alike third-party actions, migrations that reach the default branch only through a merge, and repositories that mixed v2 and v3.
\todo{Independent validation on the sample (manual\_validation.xlsx): A -- of 30 randomly sampled migrations, X removed the deprecated use as detected (precision); C -- of the 14 non-migrated pairs, Y still contained an active deprecated use in the collected snapshot; B -- agreement of the author classification with GitHub's author-association badge on 30 migrations.}

\section{Data Analysis}\label{sec:analysis}

\textbf{Time and censoring.}
For every pair we measure time in days from the event's deprecation notice.
Pairs that never migrated are right-censored at the repository's last commit or at data collection, whichever is earlier: a project that stopped committing had no further chance to migrate.

\textbf{RQ1.}
We assign each migration to a phase of its event's timeline.
For the artifact actions the phases are \emph{before any breakage} (before the first brownout for v3, before the cutoff for v1/v2), \emph{during the brownouts} (9--29~January~2025), and \emph{after the cutoff}.
For \so{} they are \emph{before the announced cutoff} (31~May~2023), \emph{between the announced cutoff and its postponement} (24~July~2023), and \emph{after the postponement}.
We also plot migrations per calendar week against the key dates.

\textbf{RQ2.}
We estimate the share of exposed repositories still using the deprecated component over time with the Kaplan--Meier estimator~\cite{kaplan1958} and compare the two events with the log-rank test~\cite{mantel1966}.
Because the events differ in period and in the projects they reach, we also compare, within the repositories that migrated away from both, the time each migration took (Wilcoxon signed-rank test).
As a sensitivity check we repeat the comparison with all non-migrants censored at the collection date.

\textbf{RQ3.}
We call a migration \emph{timely} if it happened before anything broke (before any breakage for the artifact actions, before the announced cutoff for \so{}).
For each pair we measured, on the default branch at $T_{warn}$: whether Dependabot was configured to update GitHub Actions and whether a Renovate configuration existed; the commits and distinct authors in the preceding six months; the number of workflow files; and the number of deprecated uses.
Stars, contributors and age come from the SEART export.
We compare timely and other pairs with the Mann--Whitney U test and Fisher's exact test, and fit, per event, a Cox proportional-hazards model~\cite{cox1972} of the time to migration with four covariates (stars, prior commits and deprecated uses on a log scale, and whether Dependabot or Renovate was configured), about one covariate per ten migrations.
We correct each family of p-values with Holm's method~\cite{holm1979} and check the proportional-hazards assumption by correlating each covariate's Schoenfeld residuals with event time~\cite{grambsch1994}.

\textbf{RQ4.}
For each migration we identify the commit that removed the last deprecated use; when the migration reached the default branch through a merge, this is the commit inside the merge, not the merge itself.
Its author is a \emph{bot} if the name identifies a dependency or automation bot (Dependabot, Renovate, GitHub Actions).
A human author is a \emph{core maintainer} if, in the year before the migration, they authored at least 10 commits or at least 10\% of the default-branch commits (and at least 3), or authored a merge commit, or if the change was merged from a branch of the repository itself, which requires write access.
Other authors are \emph{external contributors}.
We also record whether a bot initiated a change that a human finished.

\textbf{RQ5.}
For every pair we collected the commit messages of the change and of the merge that brought it in, the workflow diff, follow-up commits to the same workflows within 30 days, and, for pairs that did not migrate, the remaining uses.
We coded each pair along three dimensions (Table~\ref{tab:codebook}): the \emph{change type}, the \emph{trigger} stated in the messages, and the \emph{problems} visible in the material.
The codes were derived inductively on the hand-picked set and then applied to all 147 pairs.
\todo{Reliability of the RQ5 coding: who coded, and agreement (Cohen's kappa~\cite{cohen1960}) on the independently coded subset D of manual\_validation.xlsx.}

\begin{table}[t]
\caption{RQ5 codebook (abridged)}\label{tab:codebook}
\centering\footnotesize
\begin{tabular}{@{}lL{0.64\columnwidth}@{}}
\toprule
\multicolumn{2}{@{}l}{\textbf{Change type} (one per pair)} \\
\texttt{bump} & one-to-one replacement (new version, \texttt{\$GITHUB\_OUTPUT}) \\
\texttt{adapt} & replacement plus changes the new behaviour required \\
\texttt{remove} & step or workflow deleted instead of replaced \\
\texttt{replace} & a different mechanism takes over \\
\midrule
\multicolumn{2}{@{}l}{\textbf{Trigger} (first that applies)} \\
\texttt{bot} & authored or initiated by a dependency bot \\
\texttt{breakage} & message reports a failing build \\
\texttt{deprecation} & message names the deprecation or its warning \\
\texttt{targeted} & change aimed at the component, no reason given \\
\texttt{routine} & general update of actions or dependencies \\
\texttt{incidental} & part of unrelated work \\
\midrule
\multicolumn{2}{@{}l}{\textbf{Problems} (one or more)} \\
\texttt{incompatibility} & new version forced extra changes or a repair \\
\texttt{environment} & runner or container could not run the new version \\
\texttt{regression} & migration introduced an error fixed later \\
\texttt{multiline} & multi-line output values broke \\
\texttt{deferral} & explicit decision to keep the old version \\
\bottomrule
\end{tabular}
\end{table}
